\documentclass{ximera}
\title{HW 3}

\begin{document}

    \begin{abstract}  Due May 10.  \end{abstract}
    \maketitle


\begin{exercise}
    
     Let $G$ be a Lie group and $v \in T_eG$. Recall that:
    \begin{itemize}
        \item There is a unique left-invariant vector field $V \in \mathfrak{X}(G)$ with $V(e) = v$.
        \item The flow $\Phi^V : G \times \mathbb{R} \to G $ is defined everywhere for all time.
    \end{itemize} 
    Let  $\gamma : \mathbb{R} \to G$ be the flow-line of $V$ with $\gamma (0) = e$.
    \begin{itemize}
        \item Show that for each $g \in G$, one has 
        \[    \Phi_t^V (g) = g \cdot \gamma (t)  .       \]
        \item Show that  $\gamma $ is a group homomorphism.
    \end{itemize}

\end{exercise}
\begin{exercise}
      Let $V$ be a vector space with basis $\{ e_1, \ldots , e_m \}$ and $W$ a vector space with basis $\{ f_1, \ldots , f_n \}$. Using the following outline, show that the set
    \[  B : =  \{ e_i \otimes f_j \, \vert \, 1 \leq i \leq m , 1 \leq j \leq n \, \}                         \]
    is a basis of $V \otimes W$. 
    \begin{itemize}
        \item Show that the set $j (V \times W) \subset V \otimes W$ in the definition of tensor product spans $V \otimes W$. 
        \item Show that $B$ spans $V \otimes W$.
        \item Show that the dimension of $V \otimes W$ is at least the dimension of the vector space of bilinear maps $V \times W \to \mathbb{R}$. 
        \item Conclude using the fact that the space of bilinear maps $V \times W \to \mathbb{R}$ has dimension 
        \[   \dim (V) \dim (W)  .\]
    \end{itemize}

\end{exercise}
   

\begin{exercise}
     Let $k \geq 2$ and $V, X$  real vector spaces. A $k$-multilinear map is a map 
    \[ \phi : V \times \cdots \times V \to X \] 
    that when all but one coordinates are fixed, it depends linearly on the remaining coordinate. Namely, for each $v_1, \ldots , v_{k-1} \in V$, the maps 
    \begin{align*}
        \phi ( \cdot , v_1, \ldots , v_{k-1}) & : V \to X \\
        \phi ( v_1,  \cdot , \ldots , v_{k-1}) & : V \to X \\
        \ldots & \\
        \phi (  v_1, \ldots , v_{k-1}, \cdot ) & : V \to X  
    \end{align*}
     are linear.  The $k$-th iterated tensor product is defined inductively  as $V^{\otimes 0} = \mathbb{R}$,  $V^{\otimes 1} = V$, and
    \[    V^{\otimes k} : =  V^{\otimes (k-1)}   \otimes V .                     \]
    Show there is a $k$-multilinear function $j_k : V\times \cdots \times V \to V^{\otimes k}$ such that the following holds:
    \begin{center}
        For each $k$-multilinear map $\phi : V \times \cdots \times V \to X$, there is a unique linear map $\tilde{\phi} : V ^{\otimes k} \to X$ with $ \tilde{\phi} \circ j _k = \phi $. 
    \end{center}

\begin{hint}
          Take for granted the case $k = 2$ and use it as the base of induction.  

\end{hint}


\end{exercise}
    

    We denote $j_k (v_1, \ldots , v_k) $ by $v_1 \otimes \cdots \otimes v_k$. The map $j_{k + \ell } $ defines a product 
    \[ V^{\otimes k} \times V^{\otimes \ell} \to V^{\otimes (k + \ell )} , \]
    making $\mathcal{T}(V)$ a (non-commutative) algebra. 


    \begin{exercise}
        Let $k \geq 2$ and $V$  a real vector space with basis $\{ e_1, \ldots , e_n \}$. Show that the set 
    \[    B_k : = \{ e_{i_1} \otimes \cdots \otimes e_{i_k} \, \vert \, i_1, \ldots , i_k \in \{ 1, \ldots , n \} \,\}                          \]
    is a basis of $V ^{ \otimes k }$. 

    \end{exercise}
    

    
\begin{exercise}
     Let $V$ be a  vector space with basis $\{ e_1 , \ldots , e_n \}$. Recall that the exterior algebra $\wedge ^{\ast} V$ is the quotient of the tensor algebra $\mathcal{T}(V)$ over the ideal generated by the set 
    \[  I_{ext} ^2 : =  \text{span}( \{  a \otimes b + b \otimes a \, \vert \, a,b \in V \} )  .   \]
    For $k \in \mathbb{N}$, we denote by $\wedge ^k V$ the image of $V^{\otimes k}$ in  $\wedge ^{\ast} V$. For  $v_1 \otimes \cdots \otimes v_k \in V^{\otimes k}$, denote by $v_1 \wedge \ldots \wedge v_k \in \wedge ^{k}V$ its image in the quotient. Show that the set
    \[     B^{\wedge}_k : = \{  e_{i_1} \wedge \ldots \wedge e_{i_k}\, \vert \, 1 \leq i_1 < \ldots  < i_k \leq n  \, \}        \]
    is a basis of $\wedge ^k V$. 

    \textbf{Hint: }   A $k$-multilinear map  $ \phi : V \times \cdots \times V \to X$ is said to be \emph{alternating} if for any  vectors $v_1, \ldots , v_k \in V$ and any permutation $\sigma : \{ 1, \ldots, k \} \to \{ 1, \ldots , k \}$, one has 
    \[   \phi (v_{\sigma(1)}, \ldots , v_{\sigma (k)}) = \text{sign} (\sigma ) \,  \phi ( v_1,  \ldots , v_k  ) .     \]
     where sign$(\sigma ) \in \{ -1 , 1 \} $ denotes the sign of the permutation. Show that for any alternating $k$-multilinear map $\phi : V\times \cdots \times  V \to X$, there is a unique linear map $\overline{\phi} : \wedge ^k V \to X$ with
     \[    \phi (v_1, \ldots , v_k) = \overline{\phi} (v_1 \wedge \ldots \wedge v_k)          \]
     for all $v_1, \ldots , v_k \in V$.  First show that $\wedge ^k V = V^{\otimes k} / I_{ext}^k$ with 
    \[       I_{ext}^k =  \text{span} \left( \, \, \bigcup_{i = 0}^{k-2}  V^{\otimes i} \otimes I_{ext}^2 \otimes V^{k-2-i}   \right)   .    \]

\end{exercise}
    
     



\end{document}